% !TEX root = ../main.tex
\chapter{Atalhos de teclado}
\label{ap:atalhos}

Padrões de fábrica, todos reconfiguráveis nas Configurações, na seção Atalhos de Teclado.

\begin{longtable}{lp{0.52\linewidth}}
\toprule
Atalho & Ação \\
\midrule
\endhead
\tecla{Ctrl+N} & Novo arquivo \\
\tecla{Ctrl+S} & Salvar \\
\tecla{Ctrl+Shift+S} & Salvar todos \\
\tecla{Ctrl+W} & Fechar aba \\
\tecla{Ctrl+Shift+T} & Reabrir última aba fechada \\
\tecla{Ctrl+Shift+F} & Buscar nos arquivos do projeto \\
\tecla{Ctrl+Shift+P} & Paleta de comandos \\
\tecla{Ctrl+Alt+S} & Ligar ou desligar a análise semântica (slang) \\
\tecla{Shift+Alt+F} & Formatar documento (clang-format) \\
\bottomrule
\end{longtable}

Dentro do editor valem ainda os atalhos usuais do Monaco, herdados do VS Code: \tecla{Ctrl+F} para localizar, \tecla{Ctrl+H} para substituir, \tecla{Ctrl+Z} e \tecla{Ctrl+Y} para desfazer e refazer, \tecla{Alt+Click} para múltiplos cursores e \tecla{Ctrl+/} para comentar a linha, entre outros.
